\documentclass[12pt,a4paper]{article}
\usepackage[utf8]{inputenc}
\usepackage[russian]{babel}
\usepackage{geometry}
\geometry{margin=2cm}

\title{Инвариантное задание 12\\
Личная эффективность (на примере успешных людей из IT-сферы)}
\author{Ибрагимов Р.А.}
\date{}

\begin{document}

\maketitle

\section{Введение}

Что делает человека по-настоящему эффективным? Анализ жизни выдающихся людей — предпринимателей, инженеров и новаторов — показывает, что их успех основан не только на таланте, но и на определённых привычках, мышлении и дисциплине.

В данной работе рассматриваются принципы личной эффективности на основе примеров известных представителей IT-сферы.

\section{Примеры успешных людей}

\subsection{Стив Джобс}

Стив Джобс — сооснователь компании Apple, оказавший огромное влияние на развитие технологий и дизайна.

Основные принципы:
\begin{itemize}
    \item Фокус важнее многозадачности — умение отказываться от лишнего
    \item Внимание к деталям и качеству продукта
    \item Ориентация на будущее потребностей пользователей
    \item Осознанность в выборе жизненного пути
\end{itemize}

\subsection{Илон Маск}

Илон Маск — предприниматель и инженер, основатель компаний Tesla и SpaceX.

Ключевые подходы:
\begin{itemize}
    \item Мышление первыми принципами
    \item Высокие стандарты к себе и окружающим
    \item Готовность к риску и экспериментам
    \item Интенсивная работа над значимыми целями
\end{itemize}

\section{Анализ личной эффективности}

Рассмотренные примеры показывают, что успешные люди обладают рядом общих характеристик:

\begin{itemize}
    \item Способность концентрироваться на главном
    \item Системный подход к работе
    \item Готовность к постоянному обучению
    \item Умение принимать ошибки как часть процесса
\end{itemize}

\section{Правила личной эффективности}

На основе анализа можно выделить универсальные правила:

\begin{enumerate}
    \item Ставить конкретные цели
    \item Фокусироваться на ключевых задачах
    \item Планировать день заранее
    \item Воспринимать неудачи как опыт
    \item Развивать самодисциплину
    \item Постоянно учиться
    \item Заботиться о здоровье
    \item Окружать себя сильными людьми
    \item Делать важное, а не срочное
    \item Действовать, а не откладывать
\end{enumerate}

\section{Заключение}

Личная эффективность — это результат системной работы над собой. Примеры Стива Джобса и Илона Маска показывают, что ключ к успеху заключается в дисциплине, фокусе и стремлении к развитию.

Следование данным принципам позволяет повысить продуктивность и достигать поставленных целей.

\end{document}